\SourceLeaf{232}%
\Needspace{9\baselineskip}
\begin{center}\textsc{Remarks.}\end{center}
\begingroup\small
As early as June 1831 \textsc{Wolfgang Bolyai} had sent his son \textsc{Johann}'s »little work« to \textsc{Gauss}, but it did not reach him because of the cholera epidemic (letters from \textsc{Bolyai} to \textsc{Gauss} dated 20~June 1831 and 16~January 1832). The title of the copy in \textsc{Gauss}'s papers is in \textsc{Wolfgang Bolyai}'s own hand and reads as follows.
\begin{center}Appendix prima.\end{center}
\noindent Scientia Spatii, a veritate aut falsitate Axiomatis XI\textsuperscript{mi} Euclidei (a priori haud unquam decidenda) independens: atque ad casum falsitatis quadratura circuli geometrica.

Auctore, Auctoris filio \textsc{Johanne Bolyai}, de eadem, Geometrarum in Exercitu Caesareo Regio Austriaco Castrensium Locumtenente Primario.

De eadem means de Bolya. Bolya was the ancestral estate of the Bolyai family, near Maros Vásárhely (see also p.~235, l.~18 and the footnote).

\textsc{Johann Bolyai}'s classic treatise of just 26 pages appeared as an appendix to the first volume of his father's work.

Tentamen juventutem studiosam in elementa matheseos purae, elementaris ac sublimioris, methodo intuitiva, evidentiaque huic propria, introducendi. Maros Vásárhelyini 1832.

\noindent It bears a somewhat different title there (see the letter from \textsc{Gauss} to \textsc{Gerling} of 4~February 1844, p.~235 of this volume).

By way of clarification, \textsc{Johann Bolyai} uses the sign $\TripleParallel$ for \emph{parallel} in the sense of non-Euclidean geometry, and denotes the system of Euclidean geometry by $\Sigma$ and the system of his absolute geometry by $S$.

The original of \textsc{Gauss}'s letter of 6~March 1832, which \textsc{Wolfgang} had given to his son \textsc{Johann}, has been lost. \textsc{Gauss}'s papers contain only a copy made by \textsc{Johann}. His father sent it on 26~August 1856 to \textsc{Sartorius von Waltershausen}, who added it to \textsc{Gauss}'s papers.

On p.~222, ll.~18—19 the copy reads \emph{wo folglich} in place of \emph{und folglich}, and on p.~224, l.~18 it reads $\Sigma$ \emph{in} $S$ in place of $\Sigma$ \emph{und} $S$. Since \textsc{Gauss} often abbreviates \emph{und} with an indistinct $u\partial$, a copying error is very likely. The underlined signs $\underline{f\varphi}$ and $\underline{f(180^{0}-\varphi)}$ on p.~222 and $\underline{1}$, $\underline{2}$, $\underline{3}$, $\underline{\alpha}$, $\underline{\beta}$, $\underline{\gamma}$, $\underline{A}$, $\underline{B}$, $\underline{C}$ on p.~223 were written in red ink in the original and reproduced in pencil in the copy.
\begin{flushright}\textsc{Stäckel.}\end{flushright}
\par\endgroup
